\section{评估样例库与排障手册：把质量控制写成文档}

项目上线后，评估不能停留在一次性 benchmark。线上失败、对抗样例和历史事故要持续回灌为版本化资产，本节因此把课程的评估优先原则整理成样例库、归因流程和上线检查表，使每次策略修改都能回答“改善了什么、破坏了什么、增加了多少成本”。

\subsection{评估样例库的分层构造}
为了避免 ``新策略提升 A 指标却破坏 B 能力''，建议将评估集划分为核心样例库、对抗样例库和回归样例库三层。核心库负责覆盖主业务， 对抗库负责暴露脆弱点，回归库负责监控历史问题是否复发。

\begin{knowledgebox}{建议的样例库配比}
\begin{itemize}
    \item 核心样例库：60\%（覆盖高频真实任务）
    \item 对抗样例库：25\%（越权、提示注入、长链路干扰）
    \item 回归样例库：15\%（历史事故与线上坏例）
\end{itemize}
\end{knowledgebox}

\subsection{排障优先级：先判定是评估问题还是策略问题}
很多团队排障会直接改 prompt 或模型参数，但课程强调评估的重要性后，应先判断 ``是否因为评估器偏差导致假失败/假成功''。若评估器本身不稳，后续所有优化都可能偏航。

\begin{warningbox}{排障反模式}
\begin{enumerate}
    \item 不复现实验环境，直接根据线上截图改策略。
    \item 只看平均分，不看失败分布和失败类型迁移。
    \item 发现退化后立即回滚全部改动，导致无法定位根因。
\end{enumerate}
\end{warningbox}

\subsection{标准化排障流程}
样例库只能说明哪里失败，标准化流程负责把失败转成可验证修复。本小节按定位、归因、最小修复和全量验证四步推进，要求每一步保留输入、环境、指标与成本证据，避免看到坏例后直接反复改 prompt。
\begin{table}[H]
\centering
\begin{tabular}{p{2.5cm}p{4.2cm}p{4.9cm}}
\toprule
\textbf{步骤} & \textbf{执行动作} & \textbf{产出} \\
\midrule
定位 & 在固定评估集复现实验差异 & 失败样例清单 + 指标对比 \\
归因 & 区分评估漂移/数据偏移/策略回归 & 根因标签与证据链 \\
修复 & 最小改动验证假设 & 修复候选版本与风险说明 \\
验证 & 全量回归 + 对抗测试 + 成本复核 & 上线建议与回滚条件 \\
\bottomrule
\end{tabular}
\caption{可直接执行的 Agent 排障工作流}
\end{table}

\subsection{长表：常见故障与处理策略}
完成根因分类后，团队还需要把高频故障固化成共享语言。下面的长表把工具、上下文、重试、评估、权限、记忆和环境问题分别列出，使线上事件能够快速路由到对应 owner，并沉淀为下一版回归样例。
\begin{longtable}{p{2.7cm}p{4.0cm}p{4.9cm}}
\toprule
\textbf{故障类型} & \textbf{症状} & \textbf{处理建议} \\
\midrule
\endfirsthead
\toprule
\textbf{故障类型} & \textbf{症状} & \textbf{处理建议} \\
\midrule
\endhead
工具误调用 & 参数格式正确但业务语义错误 & 在评估中加入语义约束检查；增加工具前置确认步骤 \\
上下文漂移 & 多轮后目标偏离原任务 & 引入阶段性目标重述机制；加入计划一致性奖励 \\
重试风暴 & 失败后连续无效重试导致成本飙升 & 设置失败路由与重试上限；按错误类型分流处理 \\
评估振荡 & 小改动导致指标大幅波动 & 扩大评估样本；增加置信区间与分桶统计 \\
开集退化 & 闭集提升但开集体验变差 & 增加 open-ended 评审权重；补充真实任务回放 \\
权限越界 & 触发未经授权的写操作 & 强化动作白名单；高风险工具强制人审 \\
记忆污染 & 错误历史进入长期记忆 & 增加记忆有效期与可信度标签；支持记忆回收 \\
延迟失控 & 复杂任务响应时间不可接受 & 采用阶段式答复；并行化可独立子任务 \\
解释失真 & 给出看似合理但与行为不符的理由 & 将解释与实际执行日志绑定校验 \\
环境脆弱性 & 外部 API 抖动即导致任务失败 & 加入容错层、缓存层和降级策略 \\
\bottomrule
\end{longtable}

\subsection{验收前检查清单}
修复在离线样例上通过，并不意味着已经可以上线。本小节把课程中的评估、成本与系统约束收束成十项门槛，重点检查版本可复现、危险动作门控、失败恢复、实时观测和回滚路径是否真实可用。
\begin{importantbox}{上线前 10 项必检}
\begin{enumerate}
    \item 评估集是否版本锁定并可复现。
    \item 是否同时通过闭集与开集基线。
    \item 高风险工具是否有权限门控。
    \item 失败恢复路径是否经过压测。
    \item 推理成本是否在预算区间。
    \item 异常日志是否可追踪到具体动作。
    \item 是否定义了上线后回滚阈值。
    \item 关键指标是否有实时告警。
    \item 人工接管流程是否可用。
    \item 线上样例是否可回灌训练与评估。
\end{enumerate}
\end{importantbox}

\subsection{本章小结}
把评估与排障流程文档化，能够显著减少 ``经验主义调参''。课程强调的评估优先原则，在工程上可以直接落实为样例库、排障模板和上线检查清单。


